% Part: lambda-calculus
% Chapter: introduction
% Section: basic-pr-lambda

\documentclass[../../../include/open-logic-section]{subfiles}

\begin{document}

\olfileid{lam}{int}{bas}
\olsection{توابع بازگشتی اولیهٔ پایه \usetoken{S}{lambda definable} هستند}

\begin{lem}
توابع~$\Zero$، $\Succ$ و~$\Proj{n}{i}$ !!{lambda definable} هستند.
\end{lem}

\begin{proof}
$\Zero$ همان
$\lambd[x][\lambd[y][y]]$ است.

تابع جانشین~$\Succ$ با
$\fn{Succ}(u) = \lambd[x][\lambd[y][x(uxy)]]$ تعریف می‌شود. بیندیشید
که چرا این تعریف کار می‌کند: برای هر عددنمای~$\num{n}$ که به‌صورت یک
تکرارگر در نظر گرفته شود و هر تابع~$f$، عبارت~$\fn{Succ}(\num{n},f)$
تابعی است که با ورودی~$y$، تابع~$f$ را~$n$ بار، با آغاز از~$y$، اعمال
می‌کند و سپس یک بار دیگر آن را اعمال می‌کند.

دربارهٔ افکنش‌ها چیزی برای گفتن نیست:
$\fn{Proj}^n_i(x_0, \dots, x_{n-1}) = x_i$. به بیان دیگر، بنا بر قراردادهای
ما، $\fn{Proj}^n_i$ همان ترم لامبدای
$\lambd[x_0][\dots \lambd[x_{n-1}][x_i]]$ است.
\end{proof}

\end{document}
